% !TEX root = ../Tese_LucasMaio.tex

\chapter*{Resumo}

A Tomografia por Impedância Elétrica (TIE) é uma técnica de imageamento funcional não invasiva, portátil e livre de radiação ionizante, cuja reconstrução depende de métodos numéricos e de escolhas de visualização que influenciam diretamente a interpretação das imagens. Neste trabalho, foi desenvolvida uma interface gráfica interativa, em Python, voltada à reconstrução e visualização de dados previamente registrados de TIE, com foco educacional e de apoio à pesquisa. A aplicação integrou o framework \textit{open-source} pyEIT a uma arquitetura modular inspirada no padrão MVC (\textit{Model--View--Controller}) e incorporou três métodos de reconstrução --- Back-Projection (BP), Jacobiano (JAC) e GREIT --- além de duas implementações de interface, baseadas em Matplotlib e PyQtGraph.

Os resultados foram obtidos a partir de um conjunto experimental composto por 119 \textit{frames}, registrados em um arranjo simplificado com solução aquosa e oito eletrodos. A comparação qualitativa entre os métodos mostrou que os três foram capazes de localizar a haste na mesma região do domínio, embora com diferenças marcantes em contraste, definição espacial e estabilidade entre \textit{frames}. Também foram explorados parâmetros configuráveis da interface, como densidade da malha e suavização temporal, evidenciando seu impacto sobre a aparência da reconstrução e sobre o comportamento visual da aplicação. A interface alternativa em PyQtGraph demonstrou a viabilidade de uma segunda interface de visualização, ainda que com menor completude em relação à interface principal.

Conclui-se que o trabalho contribuiu com uma ferramenta computacional modular para exploração interativa de reconstruções em TIE, permitindo comparar algoritmos e parâmetros de forma visual e integrada. Como continuidade, permanecem relevantes a integração com sistemas reais de aquisição de dados, a adoção de métricas quantitativas de avaliação e a ampliação da interface alternativa baseada em PyQtGraph.

\textbf{Palavras-chave}: Tomografia por Impedância Elétrica; pyEIT; PyQt; reconstrução de imagens; visualização científica.